%% ma: L10-B04-0001
%% lop: 10
%% chuong: 2
%% bai: 4
%% dang: 10.4.1
%% loai: TLN
%% muc: VD
%% dap_an: "7"
%% nguon: "(NBV)-ĐỀ SỐ 4-MỖI NGÀY 1 ĐỀ THI 2025.pdf (D04, MỖI NGÀY 1 ĐỀ - Đề số 4), Phần III Câu 2"
%% kien_thuc: "Hệ phương trình bậc nhất ba ẩn; định luật Ôm (Vật lí 11)"
%% trang_thai: cach-ly
%% ly_do: "Câu thuộc môn Vật lí (định luật Ôm, mạch nối tiếp và song song), không nằm trong SGK Toán Kết nối tri thức; phần Toán chỉ là giải hệ ba phương trình bậc nhất ba ẩn. Gợi ý: nếu muốn giữ, giáo viên có thể duyệt chuyển vào kho liên môn hoặc viết lại thành bài toán hệ phương trình thuần Toán."
%% ghi_chu: "Đề gốc có hình mạch điện (không vẽ lại vì câu bị cách ly). Đáp án gốc 7 đúng (sympy). Mã dạng 10.4.1 chỉ là mã tạm để câu có nhãn hợp lệ."
\begin{ex}
Cho mạch điện như hình vẽ: điện trở $R_1$ mắc nối tiếp với hai nhánh song song gồm $R_2$ và $R_3$, nguồn có hiệu điện thế $U$ đặt vào hai đầu $A,B$ của mạch. Biết $R_1=36\ \Omega$, $R_2=90\ \Omega$, $R_3=60\ \Omega$ và $U=60\ \mathrm V$. Gọi $I_1$ là cường độ dòng điện chạy qua mạch chính, $I_2$ và $I_3$ là cường độ dòng điện chạy qua hai nhánh. Khi đó $6I_1+3I_2+2I_3$ bằng bao nhiêu ampe (A)?
\shortans{7}
\loigiai{$I_1=I_2+I_3$; $90I_2=60I_3$ (hai nhánh song song cùng hiệu điện thế); $36I_1+90I_2=60$. Giải hệ: $I_1=\frac56$, $I_2=\frac13$, $I_3=\frac12$. Do đó $6I_1+3I_2+2I_3=5+1+1=7$.}
\end{ex}
